\documentclass[11pt,a4paper]{article}
\usepackage[T1]{fontenc}
\usepackage[utf8]{inputenc}
\usepackage{lmodern,amsmath,amssymb}
\usepackage[margin=25mm]{geometry}
\usepackage[hidelinks]{hyperref}
\hypersetup{pdftitle={A finite-speed return bridge and its polygonal action},pdfauthor={navstokgap research working notes}}
\newcommand{\tightlist}{\setlength{\itemsep}{0pt}\setlength{\parskip}{0pt}}
\setlength{\emergencystretch}{3em}
\setcounter{secnumdepth}{0}
\begin{document}
\hypertarget{a-finite-speed-return-bridge-and-its-polygonal-action}{%
\section{A finite-speed return bridge and its polygonal
action}\label{a-finite-speed-return-bridge-and-its-polygonal-action}}

A velocity-resolved telegraph bridge gives one consistent random
trajectory at every cut resolution. Its midpoint has a discrete
probability mass, while the action error of sampled polygons vanishes at
a controlled rate. The construction isolates velocity memory from the
introduction of fresh noise at each cut.

\hypertarget{model-and-endpoint-convention}{%
\subsection{1. Model and endpoint
convention}\label{model-and-endpoint-convention}}

Fix mass \(m>0\), speed \(0<u<c\), reversal rate \(\lambda>0\) and
duration \(T>0\). In an inertial frame set \(X_0=0\), \(V_0=u\) and

\[V_t=u(-1)^{N_t},\qquad X_t=\int_0^t V_s\,ds,\]

where \(N\) is a homogeneous Poisson process of rate \(\lambda\).
Velocity is right-continuous. The rate is a supplied classical
stochastic clock. This reduced model describes impulsive reversals;
momentum receivers require the extra interaction bookkeeping of the
physical-cut model (C030).

We construct the bridge to \((X_T,V_T)=(0,-u)\) by disintegrating
switch-time densities at the interior position \(X_T=0\). The explicit
simplex construction below fixes the version of this zero-probability
conditioning. Before conditioning, the zero-switch event has probability
\(e^{-\lambda T}\) and lies at \((uT,+u)\); the chosen terminal velocity
selects odd positive switch counts.

The free kinetic functional is \(S[X]=(m/2)\int_0^T V_t^2dt\), in action
units. We compare it with the sampled polygon and with the zero-position
chord. The chord matches position endpoints; its velocity endpoints
differ from the bridge. All comparisons below name this positional
endpoint convention.

\hypertarget{exact-conditional-path-law}{%
\subsection{2. Exact conditional path
law}\label{exact-conditional-path-law}}

Write \(N_T=2K+1\) and \(z=\lambda T\). The bridge count law is

\[\boxed{\Pr(K=k\mid\text{bridge})=
 w_k=\frac{(z/2)^{2k}}{(k!)^2 I_0(z)},\qquad k=0,1,\ldots,}\]

where \(I_0(z)=\sum_{k\ge0}(z/2)^{2k}/(k!)^2\). Given \(K=k\),
independently choose uniform simplex vectors \((P_0,\ldots,P_k)\) and
\((M_0,\ldots,M_k)\), each summing to \(T/2\). Follow velocities
\(+u,-u,+u,-u,\ldots\) for durations \(P_0,M_0,P_1,M_1,\ldots,P_k,M_k\).
For \(k=0\) each simplex is a singleton. This specifies the entire
conditioned path, including all later observations.

\textbf{Derivation.} Given \(N_T=2k+1\), the ordered switch times are
uniform in \(0<s_1<\cdots<s_{2k+1}<T\). Their \(2k+2\) spacings are
uniform on the simplex of total length \(T\). The total positive
duration \(A\) has density

\[f_k(a)=\frac{(2k+1)!}{(k!)^2 T^{2k+1}}
              a^k(T-a)^k,\qquad 0<a<T.\]

Since \(X_T=u(2A-T)\), the position Jacobian is \(da/dx=1/(2u)\).
Multiplication by the Poisson count weight gives the joint density

\[\left.\frac{\Pr(X_T\in dx,N_T=2k+1)}{dx}\right|_{x=0}
 =\frac{e^{-\lambda T}\lambda^{2k+1}}{2u(k!)^2}(T/2)^{2k}.\]

Its sum is \(e^{-\lambda T}\lambda I_0(z)/(2u)>0\). Division gives
\(w_k\). Conditioning the uniform spacings on \(A=T/2\) leaves the two
independent simplexes stated above. The occupation-time backbone is the
equal-rate specialization of Cinque's Theorem 2.1, (2.6), printed p.~4;
\href{https://arxiv.org/abs/2202.01904v1}{the source} also gives the
alternating switch representation (2.4), p.~3.

\hypertarget{midpoint-law-including-the-atom}{%
\subsection{3. Midpoint law, including the
atom}\label{midpoint-law-including-the-atom}}

The exact midpoint law is an explicit pushforward of the preceding
mixture. Let \(d=(P_0,M_0,\ldots,P_k,M_k)\), \(s_0=0\) and
\(s_j=\sum_{i<j}d_i\). Then

\[Y=X_{T/2}=u\sum_{j=0}^{2k+1}(-1)^j
        \min\{d_j,\max(0,T/2-s_j)\}.\]

For any bounded measurable \(g\), its expectation is
\(\sum_k w_k\mathbb E_k[g(Y)]\), where \(\mathbb E_k\) is integration
over the two normalized simplexes. This formula sums over midpoint
velocities. The joint law retains \(V_{T/2}=u(-1)^j\) on
\(s_j\le T/2<s_{j+1}\).

In particular,

\[\boxed{\Pr(Y=uT/2)=\frac1{I_0(\lambda T)}.}\]

For \(K=0\) the unique switch occurs at \(T/2\), and \(Y=uT/2\) with
right-continuous velocity \(-u\). For \(k\ge1\), all simplex coordinates
are strictly positive almost surely. Attaining \(Y=uT/2\) would use all
positive duration before the midpoint, leaving extra positive intervals
with zero length. Attaining \(Y=-uT/2\) would require zero initial
positive duration. Both are simplex boundary events. The midpoint lies
in an interval with index \(1\le j\le2k\). For odd \(j\) its position is
\(u(2\sum_{i<j,\ i\text{ even}}d_i-T/2)\); for even \(j\) it is
\(u(T/2-2\sum_{i<j,\ i\text{ odd}}d_i)\). Each sum is a nonempty proper
prefix of its simplex vector, so on each such region \(Y\) is a
nonconstant affine function of the free simplex coordinates. Thus the
\(k\ge1\) contribution is absolutely continuous on \((-uT/2,uT/2)\).

\textbf{The velocity convention matters at the atom.} On the one-switch
component, conditioning instead on \(|X_T|<\epsilon\) makes the switch
time uniform in a symmetric interval around \(T/2\). Half of these paths
still have velocity \(+u\) at the midpoint and half have \(-u\). As
\(\epsilon\downarrow0\), their position paths converge uniformly to the
one-switch return path, but their midpoint velocity laws retain that
equal mixture. Evaluation of velocity at a jump is discontinuous. Our
exact bridge uses the right-continuous path version above; an
endpoint-window measurement is a separate specified protocol.

\hypertarget{cuts-sample-one-preparation}{%
\subsection{4. Cuts sample one
preparation}\label{cuts-sample-one-preparation}}

Every finite set of cuts is evaluated on this same count/simplex
probability space. If \(\rho\) refines \(\pi\), deleting the extra
coordinates of \((X_t,V_t)_{t\in\rho}\) returns \((X_t,V_t)_{t\in\pi}\)
pointwise. Hence their joint laws satisfy exact restriction consistency.
This argument handles atoms without multiplying singular transition
densities.

An independent midpoint-reset rule fails a direct test. Conditional on
\(Y=uT/2\), speed at most \(u\) forces velocity \(+u\) almost everywhere
on the first half and \(-u\) almost everywhere on the second half. In
particular \(X_{T/4}=X_{3T/4}=uT/4\) deterministically. Giving either
quarter point fresh nonzero variance changes the old preparation or
violates its speed support. Velocity and the endpoint conditioning carry
the required memory.

\hypertarget{exact-action-limit-under-arbitrary-cuts}{%
\subsection{5. Exact action limit under arbitrary
cuts}\label{exact-action-limit-under-arbitrary-cuts}}

For any partition \(\pi\) of \([0,T]\), let \(X^\pi\) be the sampled
linear interpolant, \(|\pi|\) its largest interval, and
\(S_\pi=S[X^\pi]\). Then

\[S[X]=\frac{mu^2T}{2},\qquad
0\le S[X]-S_\pi\le\frac{mu^2}{2}N_T|\pi|.\]

Indeed the polygon velocity on interval \(I\) is the average
\(\bar V_I\), and

\[S[X]-S_\pi=\frac m2\sum_{I\in\pi}
                  \int_I(V_t-\bar V_I)^2dt.\]

An interval without an interior switch contributes zero. Every other
interval contributes at most \(mu^2|I|/2\); their total length is at
most \(N_T|\pi|\). The count is finite almost surely. Differentiating
the convergent positive series for \(I_0\), with \(I_1=I_0'\), also
gives

\[\mathbb E[N_T\mid\text{bridge}]
       =1+z\frac{I_1(z)}{I_0(z)}.\]

Consequently \(S_\pi\to mu^2T/2\) almost surely and in \(L^1\) on any
deterministic shrinking meshes, with the displayed expected error bound.
Also \(\|X^\pi-X\|_\infty\le2u|\pi|\). Along nested partitions \(S_\pi\)
increases, by square completion at each inserted node.

The surviving action relative to the zero-position chord is the kinetic
cost of the fixed returning path. The polygon approximation error tends
to zero. Its value \(mu^2T/2\) tends to zero with duration or speed.
These are separate limits from A01's long-observation stationary
coefficient \(H_*=mu^2/\lambda\). For the bridge midpoint observable,
C031 gives \(0\le\kappa_{\rm mid}=4m\operatorname{Var}(Y)/T\le mu^2T\).
Its mean is generally biased, so midpoint action uses
\(2m\mathbb E[Y^2]/T\), rather than variance alone.

\hypertarget{consequence-for-the-action-gap-programme}{%
\subsection{6. Consequence for the action-gap
programme}\label{consequence-for-the-action-gap-programme}}

Finite speed, velocity memory and exact cut consistency coexist with
vanishing polygon error and arbitrarily small action scales. Their
combination provides a concrete classical admissible family against
which a proposed selection principle can be tested. A positive universal
remainder requires an additional restriction on this family, traced to
its physical role.

The next test is force control: replace impulses by bounded-acceleration
turns and determine the minimum duration of a return with prescribed
opposite endpoint velocities. This introduces an energy/force timescale
through a mechanical constraint, making its action dependence and
scaling testable.

Claims and source coverage are maintained separately in the ledger and
the B14 audit. The present path law is a specialization of established
telegraph probability; the cut/action consequences are derived here for
the project test.
\end{document}
